\documentclass[12pt,a4paper,oneside]{article}
\usepackage[brazil]{babel}
\usepackage[utf8]{inputenc}
\usepackage{olymp}

\setcounter{problem}{0}

\begin{document}

\begin{problem}{Nomes dos países}{paises}{2}

A Copa do Mundo 2026 começou ontem, com vitória do México sobre a África do Sul no estádio Azteca, na Cidade do México. Nesse estádio histórico, já foram disputadas duas finais de Copa do Mundo:
o Brasil de Pelé foi campeão em 1970 e a Argentina de Maradona foi campeã em 1986.

Ele não será sede da final desta Copa porque ela é disputada também no Canadá e nos Estados Unidos, e a final será no MetLife Stadium, em Nova Jersey, onde o Brasil estréia amanhã contra o Marrocos.

À noite, em Guadalajara, Coréia do Sul venceu de virada a República Tcheca. Vale informar que há 10 anos a República Tcheca registrou o nome curto e informal Tchéquia. Embora não comum no Brasil,
é o nome recomendado para uso no dia a dia, turismo e esportes. República Tcheca continua como nome formal, reservado para documentos oficiais e governamentais. Algo semelhante a Brasil e República Federativa do Brasil.

Esta Copa é a primeira com 48 países, o que permitiu algumas seleções se classificarem pela primeira vez: Cabo Verde, Curaçao, Jordânia e Uzbequistão.
Mesmo com tantas vagas, a Itália, tetracampeã mundial, conseguiu a proeza de não se classificar novamente, ao perder a repescagem nos pênaltis para a Bósnia-Herzegovina. A última vez que a Itália disputou uma Copa foi em 2014, no Brasil.
Curiosamente, para a de 2022, ela foi eliminada na repescagem pela Macedônia do Norte, anteriormente chamada Antiga República Iugoslava da Macedônia, um país que nunca se classificou para a Copa e que foi parte da Iugoslávia, junto com a Bósnia, Croácia e outros.

Nesta questão, você precisa apenas informar a quantidade de palavras do nome de um país. Por exemplo, Brasil tem 1, República Tcheca tem 2 e Coréia do Sul tem 3.

As palavras são separadas apenas por espaços, não há hífens na entrada. Por exemplo, não precisa se preocupar com ``Bósnia-Herzegovina'' ou ``Reino Unido da Grã-Bretanha e Irlanda do Norte''.

%\Notes
%
%Observações, se tiver.

\InputFile

A entrada possui apenas uma linha, contendo o nome de um país.
O nome contém apenas letras (maiúsculas ou minúsculas) e espaços em branco.
%Ou seja, não há acentos, cedilhas, hífens, entre outros.
Não há espaços antes ou depois do nome, nem dois espaços seguidos, apenas entre as palavras.
Restrições: o nome tem no máximo $50$ caracteres.

\OutputFile

Seu programa deve gerar uma linha na saída, informando o número de palavras do nome do país

%\Notes

\Example

\begin{example}
\exmp{
Brasil
}{
1
}%
\end{example}

\begin{example}
\exmp{
Republica Tcheca
}{
2
}%
\end{example}

\begin{example}
\exmp{
Tchequia
}{
1
}%
\end{example}

\begin{example}
\exmp{
Antiga Republica Iugoslava da Macedonia
}{
5
}%
\end{example}

%Comentários sobre o exemplo, se necessário.

\end{problem}

\end{document}